\documentclass[12pt,a4paper]{article}

\usepackage{fontspec}
\usepackage{polyglossia}
\setmainlanguage{vietnamese}
\IfFontExistsTF{Times New Roman}{\setmainfont{Times New Roman}}{\setmainfont{TeX Gyre Termes}}

\usepackage[a4paper,margin=2.5cm]{geometry}
\usepackage{amsmath}
\usepackage{booktabs}
\usepackage{longtable}
\usepackage{tabularx}
\usepackage{array}
\usepackage{enumitem}
\usepackage{graphicx}
\usepackage{xcolor}
\usepackage{listings}
\usepackage{hyperref}
\usepackage{float}

\hypersetup{
  colorlinks=true,
  linkcolor=blue!50!black,
  urlcolor=blue!60!black,
  citecolor=blue!50!black
}

\lstdefinestyle{code}{
  basicstyle=\ttfamily\small,
  breaklines=true,
  frame=single,
  columns=fullflexible,
  keepspaces=true,
  showstringspaces=false,
  backgroundcolor=\color{gray!6}
}
\lstset{style=code}

\newcommand{\projectname}{Gomoku/Caro 15x15 AI}
\newcommand{\repoURL}{\url{https://github.com/Azios1010/tictactoe}}
\newcolumntype{Y}{>{\raggedright\arraybackslash}X}

\title{\textbf{\projectname}\\
\large Hệ thống AI kết hợp tìm kiếm cổ điển, tri thức threat và mô hình cố vấn học giám sát}
\author{
  Nhóm thực hiện: \textit{điền tên các thành viên}\\
  Môn học: \textit{điền tên học phần}\\
  Giảng viên: \textit{điền tên giảng viên}
}
\date{\today}

\begin{document}
\maketitle

\begin{abstract}
Báo cáo trình bày quá trình xây dựng hệ thống chơi Gomoku/Caro trên bàn cờ 15x15. Thành phần ra quyết định chính của dự án là một engine AI cổ điển sử dụng minimax, alpha-beta pruning, iterative deepening, sinh nước đi ứng viên, move ordering, threat detection, heuristic evaluation, Zobrist hashing và transposition table. Bên cạnh đó, dự án mở rộng thêm một mô hình CNN policy-value học giám sát từ dữ liệu self-play để đóng vai trò cố vấn: mô hình gợi ý các nước có xác suất cao, còn engine cổ điển vẫn chịu trách nhiệm xử lý chiến thuật bắt buộc như thắng ngay, chặn thắng và đánh giá tìm kiếm. Kết quả thực nghiệm cho thấy mô hình đạt top-1 accuracy 40.244\%, top-3 accuracy 59.974\% và top-5 accuracy 69.194\% trên tập test, đồng thời illegal top-1 sau legal mask bằng 0. Hệ thống được triển khai full-stack với frontend React/Vite, backend FastAPI và arena self-play để sinh dữ liệu. Báo cáo cũng nêu rõ giới hạn: đây chưa phải engine Gomoku cấp thi đấu, chưa tích hợp Threat Space Search/VCF solver đầy đủ và chưa chứng minh mạnh hơn các engine chuyên nghiệp như Rapfi hoặc Yixin.
\end{abstract}

\noindent\textbf{Từ khóa:} Gomoku, Caro, minimax, alpha-beta pruning, heuristic search, threat detection, CNN policy-value, policy prior, FastAPI, React.

\tableofcontents
\newpage

\section{Thông Tin Chung}

\begin{longtable}{p{0.30\textwidth}p{0.60\textwidth}}
\toprule
\textbf{Mục} & \textbf{Nội dung} \\
\midrule
Tên đề tài & Xây dựng AI cho game Gomoku/Caro 15x15 \\
Loại bài toán & Game đối kháng, turn-based, zero-sum \\
Repository & \repoURL \\
Tên thư mục & \path{Tictactoe}; tên cũ của repo, nhưng bài toán hiện tại là Gomoku/Caro 15x15 \\
Frontend & React + Vite \\
Backend & Python FastAPI + Pydantic + Uvicorn \\
AI chính & Classical game engine \\
Thuật toán chính & Minimax, alpha-beta pruning, iterative deepening \\
Tri thức miền & Threat detection, heuristic pattern evaluator, immediate win/block \\
Thành phần mở rộng & Supervised CNN consultant model và hybrid policy prior \\
\bottomrule
\end{longtable}

\section{Giới Thiệu Bài Toán}
\label{sec:problem}

Gomoku/Caro là trò chơi hai người, mỗi người lần lượt đặt quân lên bàn cờ. Trong dự án này, bàn cờ có kích thước 15x15. Một bên chiến thắng khi tạo được năm quân liên tiếp theo hàng ngang, hàng dọc hoặc hai đường chéo.

Quy ước board trong engine:

\begin{longtable}{p{0.20\textwidth}p{0.70\textwidth}}
\toprule
\textbf{Giá trị} & \textbf{Ý nghĩa} \\
\midrule
\texttt{0} & Ô trống \\
\texttt{1} & AI/O/black trong internal engine \\
\texttt{-1} & Người chơi/X/white \\
\bottomrule
\end{longtable}

Bài toán AI là chọn nước đi hợp lý trong không gian trạng thái rất lớn. Ở đầu game, bàn cờ có 225 ô trống. Nếu minimax xét toàn bộ ô trống thì branching factor tăng rất nhanh, không phù hợp cho ứng dụng tương tác. Vì vậy, dự án kết hợp search cổ điển với các kỹ thuật giảm nhánh và tri thức đặc thù của Gomoku.

\subsection{Mục tiêu}

Hệ thống hướng tới các mục tiêu sau:

\begin{enumerate}[leftmargin=*]
  \item Chọn nước đi hợp lý trong thời gian ngắn.
  \item Xử lý đúng các tình huống chiến thuật cơ bản như thắng ngay, chặn thắng, tạo hoặc chặn open-four và double-threat.
  \item Trả về thông tin giải thích gồm \texttt{reason}, \texttt{evaluation} và \texttt{completed\_depth}.
  \item Hoạt động trong ứng dụng full-stack gồm frontend, backend và chế độ arena self-play.
  \item Có dữ liệu benchmark và regression test để đánh giá tiến bộ.
  \item Mở rộng bằng supervised learning advisor mà không thay thế hoàn toàn engine cổ điển.
\end{enumerate}

Mục tiêu của dự án không phải là vượt qua các engine thi đấu như Rapfi hoặc Yixin. Trọng tâm của dự án là minh họa cách kết hợp classical search với threat knowledge để tạo một bot Gomoku thực tế, có thể giải thích và phù hợp phạm vi học phần.

\section{Kiến Trúc Hệ Thống}
\label{sec:system}

\subsection{Tổng quan luồng xử lý}

\begin{lstlisting}[language={}]
Frontend React/Vite
  -> Backend FastAPI
      -> GomokuAI classical engine
      -> Consultant CNN advisor

Frontend React/Vite
  -> Arena FastAPI service
      -> Self-play engine
      -> JSONL data generation
\end{lstlisting}

Luồng người chơi với AI:

\begin{lstlisting}[language={}]
Người chơi click một ô
-> Frontend đặt quân -1
-> Frontend gọi POST /api/get-move
-> Backend validate board và difficulty
-> GomokuAI.get_move_analysis()
-> Backend trả row, col, evaluation, reason, completed_depth
-> Frontend đặt quân AI 1 và hiển thị lý do nước đi
\end{lstlisting}

\subsection{Thành phần trong repository}

\begin{longtable}{p{0.36\textwidth}p{0.54\textwidth}}
\toprule
\textbf{Thành phần} & \textbf{Vai trò} \\
\midrule
\path{frontend/src/App.jsx} & UI chính, play mode, arena mode và consultant advisor toggle \\
\path{frontend/src/components/Board.jsx} & Render board 15x15 \\
\path{frontend/src/components/Square.jsx} & Render từng ô, quân cờ và advisor badge \\
\path{backend/main.py} & FastAPI API cho chế độ người chơi với AI \\
\path{backend/ai_core.py} & Search orchestration, minimax, alpha-beta, iterative deepening và transposition table \\
\path{backend/ai_types.py} & Constants, dataclass config và result \\
\path{backend/board_rules.py} & Luật board cơ bản, winner check, normalize board \\
\path{backend/threats.py} & Threat detection theo pattern \\
\path{backend/evaluator.py} & Heuristic evaluator \\
\path{backend/move_ordering.py} & Candidate generation, tactical score và policy prior ordering \\
\path{arena/} & Self-play engine và arena service \\
\path{dl/model.py} & CNN policy-value model \\
\path{dl/predict_policy.py} & Inference wrapper cho consultant model \\
\path{model/} & Checkpoint và metrics model đã train \\
\path{tests/} & Unit tests và tactical regression \\
\bottomrule
\end{longtable}

\subsection{API chính}

\begin{longtable}{p{0.28\textwidth}p{0.62\textwidth}}
\toprule
\textbf{Endpoint} & \textbf{Vai trò} \\
\midrule
\texttt{GET /api/health} & Kiểm tra backend đang chạy \\
\texttt{POST /api/get-move} & Sinh nước đi AI chính theo difficulty \\
\texttt{POST /api/get-consultation} & Trả về các gợi ý từ CNN consultant model \\
\texttt{POST /api/report-game-result} & Ghi nhận kết quả ván chơi để phân tích \\
\texttt{GET /arena/api/health} & Kiểm tra arena service \\
\texttt{POST /arena/api/self-play} & Chạy AI tự đấu và sinh dữ liệu JSONL \\
\bottomrule
\end{longtable}

\section{Phương Pháp AI Chính}
\label{sec:method}

\subsection{Minimax và alpha-beta pruning}

Engine chính dùng minimax để mô phỏng quá trình hai bên lần lượt chọn nước đi. AI cố gắng tối đa hóa điểm đánh giá, còn đối thủ cố gắng tối thiểu hóa điểm đó. Do Gomoku 15x15 có branching factor lớn, minimax thuần không đủ nhanh. Vì vậy, dự án dùng alpha-beta pruning để bỏ qua các nhánh không còn khả năng ảnh hưởng đến quyết định cuối cùng.

\subsection{Iterative deepening và time limit}

Thay vì cố định một độ sâu lớn, AI tìm kiếm theo từng độ sâu tăng dần. Nếu hết thời gian, hệ thống trả về nước tốt nhất đã biết ở độ sâu gần nhất hoàn thành. Cách này giúp AI ổn định trong môi trường tương tác: Easy trả nhanh, Medium cân bằng giữa chất lượng và độ trễ, Hard được phép tìm lâu hơn.

\subsection{Candidate generation}

Không xét toàn bộ 225 ô, engine chỉ sinh các nước ứng viên quanh vùng đã có quân, theo bán kính và giới hạn số ứng viên của từng difficulty. Điều này giảm mạnh branching factor. Với board trống, AI ưu tiên đánh trung tâm để có thế phát triển cân bằng.

\subsection{Threat detection và heuristic evaluator}

Gomoku phụ thuộc nhiều vào pattern chiến thuật. Engine nhận diện các threat như open-four, closed-four, open-three, broken-three và double-threat. Evaluator kết hợp điểm tấn công và phòng thủ để đánh giá board. Các tình huống bắt buộc như thắng ngay hoặc chặn đối thủ thắng được xử lý trước search sâu.

\subsection{Move ordering}

Alpha-beta hiệu quả hơn nếu các nước mạnh được xét trước. Dự án xếp hạng ứng viên bằng local shape score, threat score, khả năng tấn công/phòng thủ và, ở Medium/Hard, thêm policy prior từ CNN consultant. Policy prior chỉ là điểm cộng để sắp thứ tự, không loại bỏ nước đi hợp lệ.

\subsection{Zobrist hashing và transposition table}

Trong search, cùng một thế cờ có thể xuất hiện qua nhiều thứ tự nước đi khác nhau. Zobrist hashing tạo khóa nhanh cho board và side-to-move. Transposition table lưu kết quả đã tính để giảm lặp lại, giúp search sâu hơn trong cùng giới hạn thời gian.

\section{CNN Consultant Model và Hybrid Policy Prior}
\label{sec:consultant}

\subsection{Vai trò của model}

Mô hình học giám sát không thay thế thuật toán chính. Model đóng vai trò cố vấn policy-value:

\begin{itemize}[leftmargin=*]
  \item Policy head dự đoán xác suất cho 225 ô.
  \item Value head ước lượng giá trị trạng thái.
  \item Legal mask loại bỏ ô không hợp lệ trước khi chọn top-k.
  \item Backend có endpoint \texttt{/api/get-consultation} để hiển thị các gợi ý model.
  \item Engine Medium/Hard dùng policy prior để ưu tiên thứ tự xét nước ở root search.
\end{itemize}

Cách tích hợp này thận trọng hơn so với để model tự đánh. Nếu model gợi ý chưa tốt, engine vẫn có immediate win/block, threat detection và alpha-beta search bảo vệ các quyết định chiến thuật quan trọng.

\subsection{Kiến trúc model}

Input của model là board 15x15 được mã hóa thành nhiều kênh đặc trưng, ví dụ quân của người đang xét, quân đối thủ, ô trống hoặc lượt đi. Backbone CNN học đặc trưng không gian cục bộ của Gomoku. Hai head đầu ra gồm policy logits kích thước 225 và value scalar.

\subsection{Dữ liệu huấn luyện}

Dữ liệu được sinh từ arena/self-play và lưu dạng JSONL. Mỗi sample gồm board, phân phối policy hoặc nước mục tiêu, reward/outcome và các thông tin phụ. Khi thiết kế notebook train trên Kaggle, dữ liệu được chia theo file-level split để giảm rò rỉ giữa train/validation/test.

\begin{longtable}{p{0.32\textwidth}p{0.24\textwidth}p{0.24\textwidth}}
\toprule
\textbf{Nguồn dữ liệu} & \textbf{Số file} & \textbf{Số sample} \\
\midrule
\path{data/} & 79 & 945,506 \\
\path{data/additional/} & 70 & 836,244 \\
\textbf{Tổng} & \textbf{149} & \textbf{1,781,750} \\
\bottomrule
\end{longtable}

\subsection{Kết quả model}

\begin{longtable}{p{0.40\textwidth}p{0.40\textwidth}}
\toprule
\textbf{Metric} & \textbf{Kết quả trên test set} \\
\midrule
Test samples & 50,000 \\
Top-1 accuracy & 40.244\% \\
Top-3 accuracy & 59.974\% \\
Top-5 accuracy & 69.194\% \\
Illegal top-1 trước legal mask & 65.078\% \\
Illegal top-1 sau legal mask & 0.000\% \\
Policy loss & 2.287984 \\
Value MAE & 0.725577 \\
Kaggle CUDA latency mean & 0.982 ms \\
Kaggle CUDA latency p95 & 1.065 ms \\
\bottomrule
\end{longtable}

So với baseline random legal và center-first, top-k accuracy của model cao hơn rõ rệt. Tuy nhiên, top-1 khoảng 40\% cũng cho thấy model chưa đủ tin cậy để tự quyết định toàn bộ nước đi. Vì vậy, lựa chọn tối ưu cho báo cáo là xem model như một advisor/policy prior, còn quyết định cuối cùng vẫn do search cổ điển đảm nhận.

\subsection{Hybrid policy prior trong engine}

Trong cấu hình hiện tại, Easy không dùng policy prior để giữ độ trễ thấp. Medium và Hard dùng model tại root search:

\begin{longtable}{p{0.18\textwidth}p{0.24\textwidth}p{0.24\textwidth}p{0.24\textwidth}}
\toprule
\textbf{Difficulty} & \textbf{Vai trò} & \textbf{Policy prior weight} & \textbf{Policy prior top-k} \\
\midrule
Easy & Trả nhanh, baseline cổ điển & 0 & 0 \\
Medium & Cân bằng latency và chất lượng & 10,000 & 24 \\
Hard & Ưu tiên nước model gợi ý mạnh hơn & 20,000 & 32 \\
\bottomrule
\end{longtable}

Thuật toán tích hợp:

\begin{enumerate}[leftmargin=*]
  \item Sinh candidate cổ điển để kiểm tra tactical move.
  \item Nếu có winning move hoặc blocking win, trả ngay.
  \item Nếu không có nước bắt buộc, gọi consultant model để lấy top-k legal moves.
  \item Chuyển xác suất model thành bonus cho candidate ordering.
  \item Chạy alpha-beta/iterative deepening trên danh sách đã sắp xếp.
\end{enumerate}

Điểm quan trọng là model không hard-prune các nước còn lại. Cách này giảm rủi ro bỏ sót nước phòng thủ khi model dự đoán sai.

\section{Thực Nghiệm và Đánh Giá}
\label{sec:experiment}

\subsection{Mục tiêu đánh giá}

Thực nghiệm tập trung vào ba câu hỏi:

\begin{enumerate}[leftmargin=*]
  \item AI có xử lý đúng tactical cases cơ bản không?
  \item Model consultant có học được phân phối nước đi tốt hơn baseline không?
  \item Hybrid policy prior có thể hỗ trợ thuật toán chính mà không phá vỡ các luật chiến thuật bắt buộc không?
\end{enumerate}

\subsection{Benchmark model}

\begin{longtable}{p{0.32\textwidth}p{0.18\textwidth}p{0.18\textwidth}p{0.18\textwidth}}
\toprule
\textbf{Phương pháp} & \textbf{Top-1} & \textbf{Top-3} & \textbf{Top-5} \\
\midrule
Random legal baseline & 0.585\% & 1.714\% & 2.814\% \\
Center-first baseline & 2.352\% & 4.778\% & 6.412\% \\
CNN consultant & 40.244\% & 59.974\% & 69.194\% \\
\bottomrule
\end{longtable}

Kết quả cho thấy model học được xu hướng nước đi từ dữ liệu self-play. Dù vậy, accuracy này chưa đủ để khẳng định model luôn chọn nước tối ưu theo nghĩa chiến thuật. Trong Gomoku, một nước sai ở thế bắt buộc có thể thua ngay, nên hệ thống vẫn ưu tiên search và threat rules.

\subsection{Tactical diagnostics của model}

\begin{longtable}{p{0.28\textwidth}p{0.20\textwidth}p{0.20\textwidth}p{0.22\textwidth}}
\toprule
\textbf{Case} & \textbf{Top-1 hit} & \textbf{Top-3 hit} & \textbf{Nhận xét} \\
\midrule
AI win horizontal & Có & Có & Model đặt một trong hai đầu của open-four lên top đầu \\
Block human horizontal & Có & Có & Model nhận diện đúng ô cần chặn trong case đơn giản \\
\bottomrule
\end{longtable}

\subsection{A/B benchmark: classical search và hybrid policy prior}

Để kiểm tra model có thật sự hỗ trợ thuật toán chính hay không, nhóm chạy một A/B benchmark nhỏ trên cùng cấu hình Medium. Bản \textit{classical} tắt \texttt{policy\_prior\_weight}; bản \textit{hybrid} bật \texttt{policy\_prior\_weight=10000} và \texttt{policy\_prior\_top\_k=24}. Hai bản dùng cùng depth, candidate radius, candidate limit và time limit. Model được warm-up trước khi đo để không tính chi phí load checkpoint vào latency.

\begin{longtable}{p{0.20\textwidth}p{0.17\textwidth}p{0.12\textwidth}p{0.20\textwidth}p{0.07\textwidth}p{0.09\textwidth}}
\toprule
\textbf{Case} & \textbf{Agent} & \textbf{Move} & \textbf{Reason} & \textbf{Depth} & \textbf{Mean ms} \\
\midrule
Opening center & Classical & [7,7] & opening\_center & 0 & 0.08 \\
Opening center & Hybrid & [7,7] & opening\_center & 0 & 0.07 \\
AI win horizontal & Classical & [7,9] & winning\_move & 0 & 12.88 \\
AI win horizontal & Hybrid & [7,9] & winning\_move & 0 & 12.32 \\
Block human horizontal & Classical & [7,9] & blocking\_win & 0 & 12.61 \\
Block human horizontal & Hybrid & [7,9] & blocking\_win & 0 & 9.65 \\
Two-stones opening & Classical & [6,7] & best\_search\_score & 2 & 1247.50 \\
Two-stones opening & Hybrid & [8,7] & best\_search\_score & 2 & 1258.07 \\
Scattered midgame & Classical & [8,7] & building\_attack & 2 & 1279.52 \\
Scattered midgame & Hybrid & [8,7] & building\_attack & 2 & 1282.87 \\
Quiet midgame & Classical & [6,7] & reducing\_threat & 1 & 1327.47 \\
Quiet midgame & Hybrid & [6,7] & reducing\_threat & 1 & 1333.40 \\
\bottomrule
\end{longtable}

Kết quả A/B cho thấy hybrid policy prior không làm hỏng các quyết định bắt buộc: opening, winning move và blocking move vẫn được xử lý trước search sâu. Ở các thế không có nước thắng/chặn ngay, hybrid đạt cùng completed depth với classical search. Tuy nhiên, latency chưa giảm rõ rệt vì các case non-forcing vẫn gần chạm time limit của Medium. Điều này củng cố lựa chọn thiết kế hiện tại: model nên được trình bày như một nguồn prior hỗ trợ ordering và advisor UI, chưa nên được xem là bằng chứng AI mạnh hơn rõ rệt.

\subsection{Latency}

Trong notebook Kaggle, inference latency trung bình của model trên CUDA khoảng 0.982 ms. Khi tích hợp vào backend local CPU, đo trực tiếp cho \texttt{predict\_top\_moves} cho kết quả trung bình khoảng 1.334 ms và p95 khoảng 1.652 ms sau khi model đã được load. Đo trực tiếp handler FastAPI nội bộ khoảng 2.285 ms trung bình và p95 khoảng 2.632 ms. Lần gọi đầu tiên có thể chậm hơn do chi phí load model, vì vậy backend đã warm-up consultant model khi startup.

\subsection{Regression tests đã kiểm tra}

\begin{longtable}{p{0.42\textwidth}p{0.18\textwidth}p{0.30\textwidth}}
\toprule
\textbf{Nhóm kiểm tra} & \textbf{Kết quả} & \textbf{Ý nghĩa} \\
\midrule
Python compile cho backend, arena và dl & PASS & Các file chính không lỗi cú pháp \\
\path{tests.test_policy_prior_ordering} & PASS & Policy prior chỉ ảnh hưởng ordering khi được bật \\
\path{tests.test_tactical_cases} & PASS & Immediate win/block và case chiến thuật cơ bản còn ổn định \\
\path{tests.test_consultant_api} & PASS & API consultant, fallback model path và validation hoạt động \\
Arena smoke test & PASS & Arena sinh được sample self-play nhỏ \\
\bottomrule
\end{longtable}

\subsection{Đánh giá tác động của model lên thuật toán chính}

Model có khả năng cải thiện thuật toán chính ở khía cạnh move ordering: nếu model đưa nước tốt lên sớm, alpha-beta có thể prune tốt hơn hoặc tìm thấy biến mạnh nhanh hơn. Tuy nhiên, với top-1 khoảng 40\%, model chưa đủ mạnh để được coi là bộ quyết định độc lập. Cải thiện thực tế phụ thuộc vào từng thế cờ:

\begin{itemize}[leftmargin=*]
  \item Ở thế chiến thuật rõ ràng, engine cổ điển thường đã xử lý tốt nhờ immediate win/block và threat detection.
  \item Ở thế trung cuộc không có nước bắt buộc, model có thể giúp ưu tiên các candidate tự nhiên hơn.
  \item Nếu model học lệch từ dữ liệu self-play chưa đủ mạnh, prior có thể làm thứ tự search kém hơn; do đó weight phải được giữ ở mức vừa phải.
\end{itemize}

Vì vậy, kết luận hợp lý là: model đã bổ sung một nguồn tri thức thống kê hữu ích cho ordering và tư vấn nước đi, nhưng chưa chứng minh chắc chắn làm AI mạnh hơn trong mọi tình huống. A/B benchmark cục bộ cho thấy hybrid không phá các quyết định tactical và giữ completed depth tương đương, nhưng chưa tạo lợi thế latency rõ rệt. Để chứng minh cải thiện sức chơi, cần mở rộng A/B benchmark thành nhiều ván self-play giữa bản không prior và bản có prior, cùng time limit và difficulty.

\section{So Sánh Phương Pháp}
\label{sec:comparison}

\begin{longtable}{p{0.24\textwidth}p{0.34\textwidth}p{0.32\textwidth}}
\toprule
\textbf{Hệ thống} & \textbf{Cách tiếp cận} & \textbf{Nhận xét} \\
\midrule
Minimax cơ bản & Duyệt cây game với evaluator đơn giản & Dễ hiểu nhưng quá chậm/yếu trên board 15x15 nếu không giảm nhánh \\
Dự án này & Minimax + alpha-beta + candidate pruning + move ordering + threat detection + TT + CNN prior & Phù hợp học phần, có thể giải thích và chạy tương tác \\
Rapfi/Yixin & Engine Gomoku chuyên sâu, tối ưu nhiều tầng & Mạnh hơn đáng kể; dự án chưa so sánh trực tiếp và không tuyên bố vượt qua \\
AlphaZero-Gomoku & Self-play RL, MCTS và neural network & Hướng học tăng cường phức tạp hơn; khác mục tiêu của dự án hiện tại \\
\bottomrule
\end{longtable}

\section{Khó Khăn và Cách Xử Lý}
\label{sec:difficulties}

\begin{longtable}{p{0.32\textwidth}p{0.58\textwidth}}
\toprule
\textbf{Khó khăn} & \textbf{Cách xử lý} \\
\midrule
Branching factor rất lớn & Chỉ sinh candidate gần quân đã có, giới hạn candidate theo difficulty \\
Search dễ timeout & Dùng iterative deepening và time limit \\
AI có thể bỏ sót threat & Thêm immediate win/block, threat detection và regression tactical tests \\
Evaluator khó cân bằng tấn công/phòng thủ & Tách pattern score, local score và reason classification để dễ debug \\
Model có top-1 chưa cao & Chỉ dùng model làm advisor/prior, không hard-prune và không thay thế tactical rules \\
Lần gọi model đầu chậm & Warm-up consultant model khi backend startup \\
\bottomrule
\end{longtable}

\section{Kết Luận và Hướng Phát Triển}
\label{sec:conclusion}

Dự án đã xây dựng được một hệ thống Gomoku/Caro 15x15 hoàn chỉnh gồm frontend, backend, arena self-play và AI engine. Đóng góp chính nằm ở việc kết hợp classical AI search với tri thức miền của Gomoku: candidate generation, move ordering, threat detection, heuristic evaluation, alpha-beta pruning, iterative deepening và transposition table. Thành phần CNN consultant là phần mở rộng có giá trị báo cáo: nó cho thấy dữ liệu self-play có thể được dùng để học một policy prior, nhưng vẫn được tích hợp thận trọng để giữ tính ổn định của engine chính.

Các hướng phát triển tiếp theo:

\begin{enumerate}[leftmargin=*]
  \item Mở rộng benchmark A/B thành nhiều ván self-play để so sánh engine không prior và engine có policy prior.
  \item Mở rộng tactical benchmark suite cho broken-four, double-threat và các chuỗi forcing phức tạp.
  \item Cải thiện evaluator cho các pattern Gomoku khó hơn.
  \item Lưu best move trong transposition table để tăng hiệu quả move ordering.
  \item Nghiên cứu Threat Space Search hoặc VCF solver tối giản.
  \item Thử các kỹ thuật search nâng cao như killer move, history heuristic, PVS hoặc aspiration window.
  \item Nâng chất lượng dữ liệu train và thử kiến trúc model mạnh hơn nếu muốn tăng vai trò của neural advisor.
\end{enumerate}

\appendix

\section{Phụ Lục: Lệnh Kiểm Tra}

\begin{lstlisting}[language={}]
.\backend\venv\Scripts\python.exe -m py_compile backend\ai_types.py backend\board_rules.py backend\threats.py backend\evaluator.py backend\move_ordering.py backend\ai_core.py backend\main.py arena\engine.py arena\run_arena.py dl\model.py dl\predict_policy.py

.\backend\venv\Scripts\python.exe -m unittest tests.test_policy_prior_ordering tests.test_tactical_cases tests.test_consultant_api -v

.\backend\venv\Scripts\python.exe -m arena.run_arena --games 1 --depth 1 --candidate-radius 1 --candidate-limit 4 --max-moves 6 --no-save
\end{lstlisting}

\section{Phụ Lục: Response API Mẫu}

\begin{lstlisting}[language={}]
{
  "row": 7,
  "col": 9,
  "evaluation": 9830013,
  "reason": "blocking_win",
  "difficulty": "medium",
  "completed_depth": 0,
  "message": "Move generated successfully."
}
\end{lstlisting}

\section{Phụ Lục: Ghi Chú Nộp Bài}

Trước khi nộp, nhóm nên điền tên thành viên, tên học phần, giảng viên và link demo nếu có. Nếu biên dịch PDF, nên dùng XeLaTeX hoặc LuaLaTeX vì file dùng \texttt{fontspec} và tiếng Việt Unicode.

\end{document}
